% GENERATED FILE. Edit canonical lessons, then run npm run generate:books.
% canonical-source-hash: fnv1a64:e28266647a464871
% canonical-lessons: GU-C396-gard, GU-C396-phanas, GU-C396-sakal, GU-C396-vagvu, GU-C396-ghant

\chapter{The Guard's Lantern}
\label{ch:gu-the-guard-s-lantern}
\hlchaptermodality{396}

\begin{chapteropening}
\textbf{By the end of this chapter:} \emph{I can talk about the guard, the bell and the chain on a train.}

Complete the last lesson of chapter 396: I can talk about the guard, the bell and the chain on a train.
\end{chapteropening}

\section[\texorpdfstring{\gu{ગાર્ડ} (gārḍ)}{ગાર્ડ (gārḍ)}]{\gu{ગાર્ડ} (gārḍ) — a guard (on a train or at a gate)}

\label{lesson:GU-C396-gard}



\begin{quote}
Before the new one: say the Gujarati for to cram, then the Gujarati for to seat.
\end{quote}

\subsection*{You'll want to know: \gu{ગાર્ડ}}
\textbf{\gu{ગાર્ડ}} — \emph{gārḍ} — \textquotedblleft{}a guard (on a train or at a gate)\textquotedblright{}.

\begin{quote}
\gu{ટ્રેનના} \gu{છેલ્લા} \gu{ડબ્બામાં} \gu{ગાર્ડ} \gu{બેસે} \gu{છે}.
\end{quote}

\begin{quote}
\emph{ṭrennā chellā ḍabbāmā̃ gārḍ bese che.}
\end{quote}

\begin{quote}
The guard sits in the last carriage of the train.
\end{quote}

A masculine noun.

\begin{cousinweb}
\textbf{\gu{ગાર્ડ}} (\emph{gārḍ}) is borrowed from English \emph{guard}: on the railway, the guard who rides at the back of the train.
\end{cousinweb}

\begin{grammarlens}[title={What you've built}]
One more word for getting around.
\end{grammarlens}

\subsection*{Your turn}
\begin{itemize}
\raggedright
  \item \emph{Say it:} \emph{gārḍ}
  \item \emph{Say it:} \emph{gārḍ}, once more
  \item \emph{Say it:} \emph{ṭrennā chellā ḍabbāmā̃ gārḍ bese che.} — The guard sits in the last carriage of the train.
  \item \emph{Recall:} say the Gujarati for to cram, then the Gujarati for to seat, then say \emph{gārḍ} again
\end{itemize}

\begin{tcolorbox}[breakable,colback=teal!4,colframe=teal!35!black,title={Before you move on}]
What does \gu{ગાર્ડ} mean? (\textquotedblleft{}A guard (on a train or at a gate)\textquotedblright{}.) Say it once more.
\end{tcolorbox}
\section[\texorpdfstring{\gu{ફાનસ} (phānas)}{ફાનસ (phānas)}]{\gu{ફાનસ} (phānas) — a lantern}

\label{lesson:GU-C396-phanas}



\begin{quote}
Before the new one: say the Gujarati for to seat, then the Gujarati for a guard.
\end{quote}

\subsection*{You'll want to know: \gu{ફાનસ}}
\textbf{\gu{ફાનસ}} — \emph{phānas} — \textquotedblleft{}a lantern\textquotedblright{}.

\begin{quote}
\gu{રાત્રે} \gu{ગાર્ડ} \gu{હાથમાં} \gu{ફાનસ} \gu{રાખે} \gu{છે}.
\end{quote}

\begin{quote}
\emph{rātre gārḍ hāthmā̃ phānas rākhe che.}
\end{quote}

\begin{quote}
At night the guard keeps a lantern in his hand.
\end{quote}

A neuter noun.

\begin{cousinweb}
\textbf{\gu{ફાનસ}} (\emph{phānas}) is usually traced to Persian \emph{fānūs}, \textquotedblleft{}a lantern\textquotedblright{}, a word Persian itself took from Greek.
\end{cousinweb}

\begin{grammarlens}[title={What you've built}]
One more word for getting around.
\end{grammarlens}

\subsection*{Your turn}
\begin{itemize}
\raggedright
  \item \emph{Say it:} \emph{phānas}
  \item \emph{Say it:} \emph{phānas}, once more
  \item \emph{Say it:} \emph{rātre gārḍ hāthmā̃ phānas rākhe che.} — At night the guard keeps a lantern in his hand.
  \item \emph{Recall:} say the Gujarati for to seat, then the Gujarati for a guard, then say \emph{phānas} again
\end{itemize}

\begin{tcolorbox}[breakable,colback=teal!4,colframe=teal!35!black,title={Before you move on}]
What does \gu{ફાનસ} mean? (\textquotedblleft{}A lantern\textquotedblright{}.) Say it once more.
\end{tcolorbox}
\section[\texorpdfstring{\gu{સાંકળ} (sā̃kaḷ)}{સાંકળ (sā̃kaḷ)}]{\gu{સાંકળ} (sā̃kaḷ) — a chain}

\label{lesson:GU-C396-sakal}



\begin{quote}
Before the new one: say the Gujarati for a guard, then the Gujarati for a lantern.
\end{quote}

\subsection*{You'll want to know: \gu{સાંકળ}}
\textbf{\gu{સાંકળ}} — \emph{sā̃kaḷ} — \textquotedblleft{}a chain\textquotedblright{}.

\begin{quote}
\gu{મીરાએ} \gu{બારણાની} \gu{સાંકળ} \gu{ખોલી}.
\end{quote}

\begin{quote}
\emph{Mīrāe bārṇānī sā̃kaḷ kholī.}
\end{quote}

\begin{quote}
Mira undid the chain on the door.
\end{quote}

A feminine noun.

\begin{cousinweb}
\textbf{\gu{સાંકળ}} (\emph{sā̃kaḷ}) is a tadbhava of Sanskrit \emph{śṛṅkhala}, \textquotedblleft{}a chain\textquotedblright{}; CDIAL lists the Gujarati word (CDIAL 12580). On a train it is the chain you pull to stop it.
\end{cousinweb}

\begin{grammarlens}[title={What you've built}]
One more word for getting around.
\end{grammarlens}

\subsection*{Your turn}
\begin{itemize}
\raggedright
  \item \emph{Say it:} \emph{sā̃kaḷ}
  \item \emph{Say it:} \emph{sā̃kaḷ}, once more
  \item \emph{Say it:} \emph{Mīrāe bārṇānī sā̃kaḷ kholī.} — Mira undid the chain on the door.
  \item \emph{Recall:} say the Gujarati for a guard, then the Gujarati for a lantern, then say \emph{sā̃kaḷ} again
\end{itemize}

\begin{tcolorbox}[breakable,colback=teal!4,colframe=teal!35!black,title={Before you move on}]
What does \gu{સાંકળ} mean? (\textquotedblleft{}A chain\textquotedblright{}.) Say it once more.
\end{tcolorbox}
\section[\texorpdfstring{\gu{વાગવું} (vāgvũ)}{વાગવું (vāgvũ)}]{\gu{વાગવું} (vāgvũ) — to ring; to sound (of a bell or a drum)}

\label{lesson:GU-C396-vagvu}



\begin{quote}
Before the new one: say the Gujarati for a lantern, then the Gujarati for a chain.
\end{quote}

\subsection*{You'll want to know: \gu{વાગવું}}
\textbf{\gu{વાગવું}} — \emph{vāgvũ} — \textquotedblleft{}to ring; to sound (of a bell or a drum)\textquotedblright{}.

\begin{quote}
\gu{મંદિરમાં} \gu{સાંજે} \gu{ઢોલ} \gu{વાગે} \gu{છે}.
\end{quote}

\begin{quote}
\emph{mandirmā̃ sā̃je ḍhol vāge che.}
\end{quote}

\begin{quote}
In the temple a drum sounds in the evening.
\end{quote}

Past (it …): \gu{તે} \gu{ઢોલ} \gu{વાગ્યો}, \emph{te ḍhol vāgyo}.

\begin{cousinweb}
\textbf{\gu{વાગવું}} (\emph{vāgvũ}) is a tadbhava of Sanskrit \emph{vādyatē}, \textquotedblleft{}is made to sound\textquotedblright{} (CDIAL 11513); \gu{વગાડવું}, to play an instrument, is its causative partner.
\end{cousinweb}

\begin{grammarlens}[title={What you've built}]
One more word for getting around.
\end{grammarlens}

\subsection*{Your turn}
\begin{itemize}
\raggedright
  \item \emph{Say it:} \emph{vāgvũ}
  \item \emph{Say it:} \emph{vāgvũ}, once more
  \item \emph{Say it:} \emph{mandirmā̃ sā̃je ḍhol vāge che.} — In the temple a drum sounds in the evening.
  \item \emph{Recall:} say the Gujarati for a lantern, then the Gujarati for a chain, then say \emph{vāgvũ} again
\end{itemize}

\begin{tcolorbox}[breakable,colback=teal!4,colframe=teal!35!black,title={Before you move on}]
What does \gu{વાગવું} mean? (\textquotedblleft{}To ring; to sound (of a bell or a drum)\textquotedblright{}.) Say it once more.
\end{tcolorbox}
\section[\texorpdfstring{\gu{ઘંટ} (ghanṭ)}{ઘંટ (ghanṭ)}]{\gu{ઘંટ} (ghanṭ) — a large bell}

\label{lesson:GU-C396-ghant}



\begin{quote}
Before the new one: say the Gujarati for a chain, then the Gujarati for to ring.
\end{quote}

\subsection*{You'll want to know: \gu{ઘંટ}}
\textbf{\gu{ઘંટ}} — \emph{ghanṭ} — \textquotedblleft{}a large bell\textquotedblright{}.

\begin{quote}
\gu{સ્ટેશન} \gu{પર} \gu{ઘંટ} \gu{વાગ્યો} \gu{અને} \gu{ટ્રેન} \gu{ઊપડી}.
\end{quote}

\begin{quote}
\emph{sṭeśan par ghanṭ vāgyo ane ṭren ūpḍī.}
\end{quote}

\begin{quote}
The bell rang at the station and the train set off.
\end{quote}

A masculine noun.

\begin{cousinweb}
\textbf{\gu{ઘંટ}} (\emph{ghanṭ}) is a tadbhava of Sanskrit \emph{ghaṇṭā}, \textquotedblleft{}a bell\textquotedblright{} (CDIAL 4421): the big bell of a temple or a station. \gu{ઘંટડી} is the small one.
\end{cousinweb}

\begin{grammarlens}[title={What you've built}]
One more word for getting around.
\end{grammarlens}

\subsection*{Your turn}
\begin{itemize}
\raggedright
  \item \emph{Say it:} \emph{ghanṭ}
  \item \emph{Say it:} \emph{ghanṭ}, once more
  \item \emph{Say it:} \emph{sṭeśan par ghanṭ vāgyo ane ṭren ūpḍī.} — The bell rang at the station and the train set off.
  \item \emph{Recall:} say the Gujarati for a chain, then the Gujarati for to ring, then say \emph{ghanṭ} again
\end{itemize}

\begin{tcolorbox}[breakable,colback=teal!4,colframe=teal!35!black,title={Before you move on}]
What does \gu{ઘંટ} mean? (\textquotedblleft{}A large bell\textquotedblright{}.) Say it once more.
\end{tcolorbox}
